\section{Notation and preliminaries}\label{sec:notation}

Throughout, $\F$ is an arbitrary field and all matrices have entries in
$\F$.  We use the slicing notation standard in numerical linear
algebra~\cite{golub_matrix_2013}: $A[i:j, k:l]$ is the submatrix of $A$ with
rows $i$ through $j$ and columns $k$ through $l$; a bare colon denotes the
full range, so $A[:,1:k]$ is the block of the first $k$ columns and
$A[1:k,:]$ the block of the first $k$ rows.  We write $e_i$ for the $i$th
standard basis vector, $I_n$ for the identity, and $u^T$ for the transpose.

For an $m \by p$ matrix $X$, $\rk(X)$ is its rank and
\[
  \nl(X) = p - \rk(X)
\]
is the dimension of its \emph{right} null space.  In particular, for the
three blocks entering our criteria --- each of which has exactly $k$ columns
--- we have
\begin{gather*}
  \nl(A[1:k,1:k]) = k - \rk(A[1:k,1:k]),
  \qquad
  \nl(A[:,1:k]) = k - \rk(A[:,1:k]),\\
  \nl({A[1:k,:]}^T) = k - \rk(A[1:k,:]) .
\end{gather*}

A matrix $L$ is \emph{lower triangular} if $L[i,j] = 0$ for $i < j$, and $U$
is \emph{upper triangular} if $U[i,j] = 0$ for $i > j$.  We emphasize that
triangular factors in this paper are \emph{not} assumed invertible: zeros are
allowed on their diagonals.  A triangular matrix is invertible if and only if
its diagonal entries are all nonzero, in which case its inverse is triangular
of the same kind, and products of lower (resp.\ upper) triangular matrices
are lower (resp.\ upper) triangular.  We will also use repeatedly that if
$M$ is lower triangular and $N$ is upper triangular, then
\begin{equation}\label{eq:tri-slices}
  M[1:k,:] =
  \begin{pmatrix}
    M_k & 0
  \end{pmatrix},
  \qquad
  N[:,1:k] =
  \begin{pmatrix}
    N_k \\ 0
  \end{pmatrix},
\end{equation}
where $M_k = M[1:k,1:k]$ and $N_k = N[1:k,1:k]$; if $M$ (resp.\ $N$) is
invertible then so is $M_k$ (resp.\ $N_k$), since the leading block of a
triangular matrix carries its first $k$ diagonal entries.

We record the two elementary rank facts used throughout.

\begin{lemma}[Sylvester's rank inequality]\label{lem:sylvester}
  Let $B \in \F^{m \by k}$ and $C \in \F^{k \by p}$.  Then
  \[
    \rk(B) + \rk(C) - k \ \le\ \rk(BC) \ \le\ \min\bigl(\rk(B), \rk(C)\bigr) .
  \]
\end{lemma}

\begin{proof}
  The upper bound is standard.  For the lower bound, restrict the linear map
  $x \mapsto Bx$ to $\range(C)$: its image is $\range(BC)$ and its kernel is
  contained in $\ker(B)$, so
  $\rk(C) = \dim \range(C) \le \rk(BC) + \nl(B) = \rk(BC) + k - \rk(B)$.
\end{proof}

\begin{lemma}[invariance of rank and nullity]\label{lem:rank-invariance}
  If $X$ is an $m \by p$ matrix and $G \in \F^{m \by m}$, $H \in \F^{p \by p}$
  are invertible, then $\rk(GXH) = \rk(X)$ and $\nl(GXH) = \nl(X)$.
\end{lemma}
